\section{Model Performance Evaluation}
% TODO: include hyperparameters selection
In this section, we evaluate the performance of four models: linear
regression, linear SVR, decision tree and
LightGBM, and measure the effect of quantisation on the performance of the
model. Linear regression, linear SVR and decision tree are
implemented with \texttt{scikit-learn} \cite{scikit-learn} and
LightGBM is from \texttt{lightgbm} package.
% TODO: cite lightgbm
We first train the model on default hyperparameter then quantise the
models to see how quantisation affect the model performance. For
linear regression and linear SVR have min-max scaling and
standardisation before fitting the model to improve the accuracy.

For the decision tree and LightGBM, quantisation is simple: every
split threshold $t$ is rounded down to $\lfloor t \rfloor$, and every
leaf value is rounded to the nearest integer. Since the features are
integers, $x \le t$ holds exactly when $x \le \lfloor t \rfloor$, so
rounding the thresholds does not change any split.

Linear regression and linear SVR need fixed-point arithmetic, since
both the feature scaling and the coefficients are fractional.
To eliminate floating-point arithmetic, each fractional value is
multiplied by a large constant and rounded to an integer. The final
result is then divided by the same constant to recover the actual
value. The constant is choosen as power of two so that the division can be
reduced to bitshift
\begin{equation}
    \hat{y}(\mathbf{x}) =
    \left( c_0' + \sum_{j=1}^{m} c_j' x_j \right) \gg b,
\end{equation}
where $c_0'$ and $c_j'$ are the coefficients multiplied by $2^b$. The
exponent $b$ is chosen as large as possible such that every
intermediate value fits in 32 bits.

% TODO: reivew this part
Table~\ref{tab:model_results} gives the performance of the
four models training 11 features against the validation set extracted
from the training data and one the final reserved topology for
testing. Quantisation has a negligible effect on accuracy of all four
models. The test MAE of LightGBM and the decision tree is unchanged,
and the test MAEP of the linear models shifts by less than 0.2
percentage points, well below
the gap between the model families. Since the integer-only versions
predict as accurately as their floating-point counterparts, they can
run directly on the motes, and we carry them forward to evaluation in
simulation.

\begin{table}[htbp]
    \centering
    \caption{Prediction error of each model. MAE is on the PDR scaled
    to $[0, 65535]$, and MAEP is the MAE as a percentage of this range.}
    \label{tab:model_results}
    \begin{tabular}{l|c|c|c|c|c|c}
        \hline
        & \multicolumn{2}{c|}{Train}
        & \multicolumn{2}{c|}{Test}
        & \multicolumn{2}{c}{Test (quantised)} \\
        \cline{2-7}
        \textbf{Model}
        & MAE & MAEP (\%) & MAE & MAEP (\%) & MAE & MAEP (\%) \\
        \hline
        LightGBM          & 10638.6 & 16.233 & 13751.8 & 20.984 & 13751.7 &
        20.984 \\
        Decision Tree     & 10574.7 & 16.136 & 13765.0 & 21.004 & 13764.9 &
        21.004 \\
        Linear SVR        & 12198.8 & 18.614 & 15015.7 & 22.912 & 15116.6 &
        23.066 \\
        Linear regression & 12634.1 & 19.278 & 15300.5 & 23.347 & 15216.1 &
        23.218 \\
        \hline
    \end{tabular}
\end{table}
